\chapter{Конструкторская часть}

\section{Описание сущностей базы данных}

\renewcommand{\arraystretch}{1.3}

В соответствии с таблицей~\ref{tb:data}, содержащей данные, которые должны находиться в базе данных, можно выделить следующие таблицы:
\begin{enumerate}
	\item organizer -- таблица организаторов турниров;
	\item location -- таблица площадок (локаций) проведения матчей;
	\item team -- таблица спортивных команд;
	\item referee -- таблица судей;
	\item tournament -- таблица спортивных турниров;
	\item match -- таблица матчей турнира;
	\item player -- таблица игроков команд;
	\item coach -- таблица тренеров команд;
	\item tournament\_registration -- таблица регистраций команд на турнир;
	\item match\_participant -- таблица команд-участников матчей и их результатов;
	\item match\_referee -- таблица назначения судей на конкретные матчи.
\end{enumerate}

На основе информации о выбранной СУБД и представленной на рисунке~\ref{fig:er} диаграмме сущность-связь, мы определим для каждой из указанных таблиц структуру столбцов, их типы и установим соответствующие ограничения, которые представлены в таблицах~\ref{table:db:organizer}--\ref{table:db:match_referee}.

% ----------------------------------------------------
% Таблица 1: organizer
% ----------------------------------------------------
\begin{table}[H]
	\begin{center}
		\caption{Информация о столбцах таблицы организаторов (organizer)}
		\begin{tabular}{|c|c|>{\centering\arraybackslash}m{5cm}|>{\centering\arraybackslash}m{3.5cm}|}
			\hline
			Столбец & Тип данных & Ограничения & Значение \\
			\hline
			id & bigint & NOT NULL, PRIMARY KEY & Идентификатор организатора \\
			\hline
			company\_name & text & NOT NULL & Название компании \\
			\hline
			email & text & NOT NULL, UNIQUE, CHECK (формат email) & Электронная почта \\
			\hline
			phone & text & CHECK(формат телефона) & Номер телефона \\
			\hline
		\end{tabular}
		\label{table:db:organizer}
	\end{center}
\end{table}

% ----------------------------------------------------
% Таблица 2: location
% ----------------------------------------------------
\begin{table}[H]
	\begin{center}
		\caption{Информация о столбцах таблицы площадок (location)}
		\begin{tabular}{|c|c|>{\centering\arraybackslash}m{5cm}|>{\centering\arraybackslash}m{3.5cm}|}
			\hline
			Столбец & Тип данных & Ограничения & Значение \\
			\hline
			id & bigint & NOT NULL, PRIMARY KEY & Идентификатор площадки \\
			\hline
			name & text & NOT NULL & Название \\
			\hline
			city & text & NOT NULL & Город \\
			\hline
			capacity & bigint & CHECK (capacity $>$ 0) & Вместимость \\
			\hline
		\end{tabular}
		\label{table:db:location}
	\end{center}
\end{table}

% ----------------------------------------------------
% Таблица 3: team
% ----------------------------------------------------
\begin{table}[H]
	\begin{center}
		\caption{Информация о столбцах таблицы команд (team)}
		\begin{tabular}{|c|c|>{\centering\arraybackslash}m{5cm}|>{\centering\arraybackslash}m{3.5cm}|}
			\hline
			Столбец & Тип данных & Ограничения & Значение \\
			\hline
			id & bigint & NOT NULL, PRIMARY KEY & Идентификатор команды \\
			\hline
			name & text & NOT NULL & Название команды \\
			\hline
			country & text & NOT NULL & Страна \\
			\hline
			foundation\_date & date & NOT NULL & Дата основания \\
			\hline
		\end{tabular}
		\label{table:db:team}
	\end{center}
\end{table}

% ----------------------------------------------------
% Таблица 4: referee
% ----------------------------------------------------
\begin{table}[H]
	\begin{center}
		\caption{Информация о столбцах таблицы судей (referee)}
		\begin{tabular}{|c|c|>{\centering\arraybackslash}m{5cm}|>{\centering\arraybackslash}m{3.5cm}|}
			\hline
			Столбец & Тип данных & Ограничения & Значение \\
			\hline
			id & bigint & NOT NULL, PRIMARY KEY & Идентификатор судьи \\
			\hline
			full\_name & text & NOT NULL & ФИО судьи \\
			\hline
			category & text & NOT NULL & Категория \\
			\hline
			experience\_years & int & DEFAULT 0, CHECK ($\ge$ 0) & Опыт работы (в годах) \\
			\hline
		\end{tabular}
		\label{table:db:referee}
	\end{center}
\end{table}

% ----------------------------------------------------
% Таблица 5: tournament
% ----------------------------------------------------
\begin{table}[H]
	\begin{center}
		\caption{Информация о столбцах таблицы турниров (tournament)}
		\begin{tabular}{|c|c|>{\centering\arraybackslash}m{5cm}|>{\centering\arraybackslash}m{3.5cm}|}
			\hline
			Столбец & Тип данных & Ограничения & Значение \\
			\hline
			id & bigint & NOT NULL, PRIMARY KEY & Идентификатор турнира \\
			\hline
			organizer\_id & bigint & NOT NULL, FK (organizer) ON DELETE CASCADE & Идентификатор организатора \\
			\hline
			name & text & NOT NULL & Название турнира \\
			\hline
			sport\_type & text & NOT NULL & Вид спорта \\
			\hline
			start\_date & date & NOT NULL, CHECK ($\le$ end\_date) & Дата начала \\
			\hline
			end\_date & date & NOT NULL & Дата окончания \\
			\hline
			prize\_pool & numeric & DEFAULT 0.00, CHECK ($\ge$ 0) & Призовой фонд \\
			\hline
			status & text & DEFAULT 'registration', CHECK (status IN ...) & Статус турнира \\
			\hline
		\end{tabular}
		\label{table:db:tournament}
	\end{center}
\end{table}

% ----------------------------------------------------
% Таблица 6: match
% ----------------------------------------------------
\begin{table}[H]
	\begin{center}
		\caption{Информация о столбцах таблицы матчей (match)}
		\begin{tabular}{|c|c|>{\centering\arraybackslash}m{5cm}|>{\centering\arraybackslash}m{3.5cm}|}
			\hline
			Столбец & Тип данных & Ограничения & Значение \\
			\hline
			id & bigint & NOT NULL, PRIMARY KEY & Идентификатор матча \\
			\hline
			tournament\_id & bigint & NOT NULL, FK (tournament) ON DELETE CASCADE & Идентификатор турнира \\
			\hline
			location\_id & bigint & FK (location) ON DELETE SET NULL & Идентификатор площадки \\
			\hline
			match\_date & timestamp & NOT NULL & Дата и время проведения \\
			\hline
			stage\_name & text & NOT NULL & Стадия турнира \\
			\hline
			status & text & DEFAULT 'scheduled', CHECK (status IN ...) & Статус матча \\
			\hline
		\end{tabular}
		\label{table:db:match}
	\end{center}
\end{table}

% ----------------------------------------------------
% Таблица 7: player
% ----------------------------------------------------
\begin{table}[H]
	\begin{center}
		\caption{Информация о столбцах таблицы игроков (player)}
		\begin{tabular}{|c|c|>{\centering\arraybackslash}m{5cm}|>{\centering\arraybackslash}m{3.5cm}|}
			\hline
			Столбец & Тип данных & Ограничения & Значение \\
			\hline
			id & bigint & NOT NULL, PRIMARY KEY & Идентификатор игрока \\
			\hline
			team\_id & bigint & FK (team) ON DELETE SET NULL & Идентификатор команды \\
			\hline
			nickname & text & NOT NULL & Никнейм \\
			\hline
			real\_name & text & NOT NULL & Настоящее имя \\
			\hline
			birth\_date & date & NOT NULL & Дата рождения \\
			\hline
		\end{tabular}
		\label{table:db:player}
	\end{center}
\end{table}

% ----------------------------------------------------
% Таблица 8: coach
% ----------------------------------------------------
\begin{table}[H]
	\begin{center}
		\caption{Информация о столбцах таблицы тренеров (coach)}
		\begin{tabular}{|c|c|>{\centering\arraybackslash}m{5cm}|>{\centering\arraybackslash}m{3.5cm}|}
			\hline
			Столбец & Тип данных & Ограничения & Значение \\
			\hline
			id & bigint & NOT NULL, PRIMARY KEY & Идентификатор тренера \\
			\hline
			team\_id & bigint & FK (team) ON DELETE SET NULL & Идентификатор команды \\
			\hline
			first\_name & text & NOT NULL & Имя \\
			\hline
			last\_name & text & NOT NULL & Фамилия \\
			\hline
			experience\_years & int & DEFAULT 0, CHECK ($\ge$ 0) & Опыт работы (в годах) \\
			\hline
		\end{tabular}
		\label{table:db:coach}
	\end{center}
\end{table}

\begin{table}[H]
	\begin{center}
		\caption{Информация о столбцах таблицы судей матча (match\_referee)}
		\begin{tabular}{|c|c|>{\centering\arraybackslash}m{5cm}|>{\centering\arraybackslash}m{3.5cm}|}
			\hline
			Столбец & Тип данных & Ограничения & Значение \\
			\hline
			match\_id & bigint & NOT NULL, PK, FK (match) ON DELETE CASCADE & Идентификатор матча \\
			\hline
			referee\_id & bigint & NOT NULL, PK, FK (referee) ON DELETE CASCADE & Идентификатор судьи \\
			\hline
			role & text & NOT NULL & Роль судьи \\
			\hline
		\end{tabular}
		\label{table:db:match_referee}
	\end{center}
\end{table}

% ----------------------------------------------------
% Таблица 9: tournament_registration
% ----------------------------------------------------
\begin{table}[H]
	\begin{center}
		\caption{Информация о столбцах таблицы регистрации на турнир (tournament\_registration)}
		\begin{tabular}{|c|c|>{\centering\arraybackslash}m{5cm}|>{\centering\arraybackslash}m{3.5cm}|}
			\hline
			Столбец & Тип данных & Ограничения & Значение \\
			\hline
			tournament\_id & bigint & NOT NULL, PK, FK (tournament) ON DELETE CASCADE & Идентификатор турнира \\
			\hline
			team\_id & bigint & NOT NULL, PK, FK (team) ON DELETE CASCADE & Идентификатор команды \\
			\hline
			registration\_date & timestamp & DEFAULT current\_timestamp & Дата регистрации \\
			\hline
			status\_registration & text & DEFAULT 'pending', CHECK (status IN ...) & Статус заявки \\
			\hline
		\end{tabular}
		\label{table:db:tournament_registration}
	\end{center}
\end{table}

% ----------------------------------------------------
% Таблица 10: match_participant
% ----------------------------------------------------
\begin{table}[H]
	\begin{center}
		\caption{Информация о столбцах таблицы участников матча (match\_participant)}
		\begin{tabular}{|c|c|>{\centering\arraybackslash}m{5cm}|>{\centering\arraybackslash}m{3.5cm}|}
			\hline
			Столбец & Тип данных & Ограничения & Значение \\
			\hline
			match\_id & bigint & NOT NULL, PK, FK (match) ON DELETE CASCADE & Идентификатор матча \\
			\hline
			team\_id & bigint & NOT NULL, PK, FK (team) ON DELETE CASCADE & Идентификатор команды \\
			\hline
			score & int & DEFAULT 0, CHECK ($\ge$ 0) & Счет команды \\
			\hline
		\end{tabular}
		\label{table:db:match_participant}
	\end{center}
\end{table}

% ----------------------------------------------------
% Таблица 11: match_referee
% ----------------------------------------------------
%\begin{table}[H]
%	\begin{center}
%		\caption{Информация о столбцах таблицы судей матча (match\_referee)}
%		\begin{tabular}{|c|c|>{\centering\arraybackslash}m{5cm}|>{\centering\arraybackslash}m{3.5cm}|}
%			\hline
%			Столбец & Тип данных & Ограничения & Значение \\
%			\hline
%			match\_id & bigint & NOT NULL, PK, FK (match) ON DELETE CASCADE & Идентификатор матча \\
%			\hline
%			referee\_id & bigint & NOT NULL, PK, FK (referee) ON DELETE CASCADE & Идентификатор судьи \\
%			\hline
%			role & text & NOT NULL & Роль судьи \\
%			\hline
%		\end{tabular}
%		\label{table:db:match_referee}
%	\end{center}
%\end{table}

На рисунке~\ref{fig:erd} приведена ER-диаграмма базы данных.
	

\begin{figure}[]
	\centering
	\includegraphics[width=\textwidth]{erd1}
	\caption{Схема базы данных}
	\label{fig:erd}
\end{figure}

	
\newpage
\section{Ролевая модель}

Для реализации описанной логики на уровне базы данных (PostgreSQL) проектируется ролевая модель, включающая 3 роли базы данных. Разграничение доступа к объектам (таблицам и представлениям) выглядит следующим образом.

\begin{enumerate}
	\item Неавторизованный пользователь -- имеет права доступа только на чтение к таблицам tournament, match, team, player, location, match\_participant. Запрещены любые операции изменения данных.
	
	\item Представитель команды -- те же права, что и у неавторизованного пользователя, при этом появляются возможности для добавления, обновления и удаления данных из таблиц tournament\_registration player
	
	\item Организатор турнира -- имеет полный доступ в базе данных.
\end{enumerate}
% **1. Роль `guest_role` (Гость)**
% * **Права:** Только чтение (`SELECT`).
% * **Доступ к таблицам:** `tournament`, `match`, `team`, `player`, `location`, `match_participant` (расписание и результаты).
% * **Описание:** Запрещены любые операции изменения данных. Покрывает потребности неавторизованных пользователей в получении статистики.

% **2. Роль `team_manager_role` (Участник / Представитель команды)**
% * **Права:** Чтение (`SELECT`), добавление (`INSERT`), изменение (`UPDATE`) и удаление (`DELETE`).
% * **Доступ к таблицам:** 
% * `SELECT` — ко всем открытым данным (как у гостя).
% * `INSERT`, `UPDATE`, `DELETE` — к таблицам `tournament_registration` (управление заявками) и `player` (управление составом своей команды).
% * **Описание:** Позволяет взаимодействовать с турниром и управлять только сущностями, относящимися к конкретной команде (через проверки на стороне приложения или RLS в БД).

% **3. Роль `organizer_role` (Организатор / Администратор)**
% * **Права:** Полный CRUD-доступ (`SELECT`, `INSERT`, `UPDATE`, `DELETE`) + `EXECUTE` (запуск процедур).
% * **Доступ к таблицам:** Ко всем таблицам БД схемы `sport`.
% * **Описание:** Обладает правами на генерацию расписания (вызов хранимых процедур), занесение результатов матчей (`match_participant`), изменение статусов турниров и управление справочниками (судьи, локации, организаторы).

	
\section{Используемые триггеры и процедура}

В базе данных присутствуют два триггера и одна хранимая процедура. Один триггер активируется автоматически после завершения матча-финала турнира и автоматически завершает турнир, чей финал был окончен, что позволяет обеспечить актуальность данных. Второй триггер позволяет организаторам автоматически создавать расписание группового этапа. Хранимая процедура автоматизирует процесс регистрации команд на турнир.

\begin{figure}[h]
	\centering
	\includegraphics[width=0.50\textwidth]{trig_final}
	\caption{Схема алгоритма триггера, обновляющего статус турнира после конца матча-финала}
	\label{fig:trig_final}
\end{figure}

\begin{figure}[h]
	\centering
	\includegraphics[width=0.8\textwidth]{trig_schedule}
	\caption{Схема алгоритма триггера, создающего расписание группового этапа матчей}
	\label{fig:trig_schedule}
\end{figure}

\begin{figure}[h]
	\centering
	\includegraphics[width=\textwidth]{reg_team_proc}
	\caption{Схема алгоритма хранимой процедуры, которая осуществляет регистрацию команды на турнир}
	\label{fig:proc_reg}
\end{figure}

\clearpage
\section*{Вывод}

В данном разделе были описаны все сущности базы данных на основе
анализа предметной области, представлена ER-диаграмма сущностей, описана ролевая модель. Также был спроектирован необходимый триггер.
\clearpage
